\documentclass{article}
\usepackage[T1]{fontenc}
\usepackage{lmodern}
\usepackage{graphicx}
\usepackage{amsmath,amssymb}
\usepackage[a4paper,margin=2cm]{geometry}
\usepackage[absolute,overlay]{textpos}
\setlength{\TPHorizModule}{1mm}
\setlength{\TPVertModule}{1mm}
\usepackage[indonesian]{babel}
\usepackage{hyperref}
\setlength{\emergencystretch}{2em}
% unit: Halaman 7

\begin{document}

% === Halaman 7 (llm) ===

\section*{Spesifikasi Algoritma Rijndael}

\begin{itemize}
    \item Rijndael mendukung panjang kunci 128 bit sampai 256 bit dengan step 32 bit (yaitu 128 bit, 160, 192, ..., 256 bit).
    \item Panjang kunci dan ukuran blok dapat dipilih secara independent (tidak harus sama. Misal: panang blok 128 bit, kunci 256 bit).
    \item Setiap blok dienkripsi dalam sejumlah putaran tertentu, sebagaimana halnya pada $DES$.
    \item Karena $AES$ menetapkan panjang kunci adalah 128, 192, dan 256, maka dikenal AES-128, AES-192, dan AES-256.
\end{itemize}

\end{document}
